\section{Certificates, source binding and independent replay}
\label{sec:certificates}

\subsection{Three distinct trust layers}

An algebraically well-formed table is not automatically the census of a
particular interval relation. The implementation separates three assertions.

\paragraph{Table validity.}
\texttt{Profiles} checks nonempty positive-length atoms, nonzero nonnegative
profile vectors, positive integer multiplicities, sorted unique rows, and
all coordinate mass equations \eqref{eq:mass}. These checks establish a
valid abstract finite census. They do not authenticate the source pairings.
Boolean values are rejected where exact integers are required.

\paragraph{Source compilation.}
The native adapter invokes the maintained scalar weighted producer and
wraps its certificate with the original universe, exact atom cuts, and
packing mode. Independent verification reconstructs the radix weights,
checks the native weighted certificate against the original pairing list,
and decodes its authenticated scalar histogram using separate arithmetic
code. It returns a newly constructed trusted table. It does not return a
Boolean that a caller might mistakenly apply to an unrelated mutable table.

\paragraph{Quotient or program claim.}
Given an authenticated table, independent graph replay constructs the
quotient without calling the incremental engine. Independent grammar replay
validates backward references and computes summaries through a separate
graph-based join rather than the producer's set-union code. These checkers
establish the subsequent claim relative to that source table. They do not
turn an unauthenticated table into a knot certificate.

\subsection{A quotient certificate need not reproduce union-find decisions}

The graph in \cref{thm:quotient} depends only on the original census and
cone groups. A verifier can reconstruct it from scratch and compare the
complete claimed output census. There is no need to trust the producer's
choice of roots, compression history, activation order, or internal count
updates. Equal count totals alone are insufficient: different profiles can
have the same number of components. The checker therefore compares the
entire canonical table when a census is claimed.

A bounded literal oracle supplies a further level of independence in the
finite tests. It explicitly builds the relation on points, adds actual
full-port equivalences, and counts each component's atom intersections.
It neither uses profile-type collapse nor relies on the graph quotient
formula. The literal oracle has a point cap and is not used on huge
multiplicities.

\subsection{First-threshold certificates}

A compact claim consists of the source table, validated grammar, target $K$,
and claimed position $t$. A trusted source table may be referenced by an
already checked compilation certificate rather than recopied.

\begin{proposition}[Two adjacent prefixes certify the first threshold]
\label{prop:threshold-cert}
For $1\le t\le L(P)$, the equations
\[
 C(t-1)>K\ge C(t)
\]
certify that $t$ is the first qualifying prefix. A claim $t=0$ is checked
against the initial count. A claim that the threshold is never attained is
checked against the final count.
\end{proposition}

\begin{proof}
Every earlier prefix has count at least $C(t-1)$ by monotonicity, so none
qualifies. The second inequality shows that $t$ does. The endpoint cases
are the same monotonicity argument.
\end{proof}

The verifier reconstructs both prefixes independently using
\cref{prop:prefix}, then performs graph quotients. No list of the first
$t$ instructions is required. The bit length of $t$ is part of the
certificate size. An out-of-range prefix, forged source cut, negative
exponent, forward grammar reference, or changed cone group is rejected.

\subsection{Certifying the complete event history}

Suppose events are claimed at
$0<t_1<\cdots<t_h\le L(P)$ with before/after counts $b_i,a_i$. Check
\[
 C(t_i-1)=b_i>a_i=C(t_i)
\]
for every event. Also check that $b_1=C(0)$, that $b_{i+1}=a_i$, and that
$C(L(P))=a_h$. For an empty list, require $C(L(P))=C(0)$.

\begin{proposition}[Plateau completeness]
\label{prop:history-cert}
The preceding checks prove both correctness of every listed event and
absence of every omitted count-changing event.
\end{proposition}

\begin{proof}
The adjacent-prefix comparisons validate the listed drops. Between two
listed drops, the checked endpoints have equal counts. A nonincreasing
integer-valued sequence with equal endpoint values is constant throughout
that interval. The same argument handles the initial and final gaps.
\end{proof}

This is stronger than checking only that the final count is correct. A
producer that silently omits an interior event will violate one of the
plateau equalities. The tests delete the first and last events from sharp
examples, mutate event indices and counts, and disable producer functions
while executing the independent checker.

\subsection{What is not certified}

The full source binding reaches an interval relation, not automatically a
manifold homeomorphism type or knot diagram. Likewise, summing an old Euler
charge over a merged equivalence class gives a sum of old charges. It does
not automatically give the Euler characteristic of a topological object
formed by attaching new cells. Such cells need correctly signed ownership
charges and a separate geometric proof. Orientation, cyclic order, slope,
and pointwise attachment maps are not supplied by the profile table.

The optional native adapter also does not claim an independent proof of the
classical AHT theorem. It relies on the maintained local trace verifier for
that layer. The added checker is independent of the new compilation and
quotient producers at the layers it actually verifies. Independence is thus
an explicit dependency boundary, not the claim that no code or mathematical
lemma is shared anywhere.
